\label{sec:results}

\subsection{Experimental Details}
\paragraph{Models.}
First, we evaluate both open-weight and proprietary LMs, including Llama 3.1 (8B, 70B) and GPT-4o (\texttt{gpt-4o-2024-05-13}). In this setup, each LM generates an answer independently, without external retrieval, and provides a list of referenced paper titles. For evaluation, we verify whether the generated paper titles exist. If they do, we retrieve their corresponding abstracts to use as citations. 
For multi-paper tasks, we further evaluate other proprietary systems: Perplexity Pro,\footnote{\url{https://www.perplexity.ai/}. We used the paid subscription version for the experiments. Note that Perplexity Search does not have an API, so we use the \texttt{selenium} toolkit and save their final prediction results from the interface. Due to this, we could not retrieve their citation information.} and PaperQA2~\citep{skarlinski2024language}, a concurrent literature review agent system that uses GPT4o for reranking, summarization, and answer generation.\footnote{We use their official codebase. As PaperQA2 does not release their retrieval corpus and requires downloading PDF files offline, we downloaded PDF files of papers suggested by our retrieval pipelines and the Semantic Scholar API. Unlike PaperQA2, we do not have access to private or license-protected papers, which may limit the effectiveness of our replication to some extent.}  
Then, we evaluate models using our \textsc{OpenScholar-DataStore} (+OSDS), where we retrieve the top $N$ passages, and concatenate and feed them together with the original input.  
Lastly, we evaluate our proposed \model, leveraging our custom inference-time pipeline using trained 8B model models ({\bf OS-8B}), as well as Llama 3.1 70B and GPT4o ({\bf OS-70B}, {\bf OS-GPT4o}).  

\paragraph{Details of \model.}
We use peS2o v2 as our default datastore $\mathbf{D}$. We analyze the effect of different datastores in Appendix \ref{app_sec:pes2o_ds_analysis}. 
For $\theta_\text{bi}$ and $\theta_\text{cross}$ in \model, we use our trained bi-encoder and cross-encoder models, which consist of 110 million and 340 million parameters, respectively. We set the maximum number of papers from web search and Semantic Scholar to 10. 
For the generator LMs, we set the temperature to 0.7 and limit the maximum token count to 3,000 for response generation and 1,000 for feedback generation, and use the \texttt{vllm} package for faster inference. We train Llama 3.1 8B for two epochs on 130k training instances for two epochs, using \texttt{torchtune}.
Further details are in Appendix~\ref{app_sec:experimental_details}. 
For all models, we set the number of passages input into the generator LM to five for single-paper tasks and ten for multi-paper tasks. No few-shot demonstrations are provided, except for SciFact and PubMed, where we include one-shot demonstrations. 


\subsection{Results}
Table~\ref{table:result_multi} show scores for multiple aspects of the main baselines. 
In summary, \model achieves state-of-the performance, significantly outperforming GPT4o and their standard RAG version, as well as specialized literature review systems such as PaperQA2~\citep{skarlinski2024language} by a large margin. 


\paragraph{Single-paper tasks.} 
On single-paper tasks, \model consistently outperforms other models. OS-8B and OS-70B outperforms Llama 3.1 8B and 70B with and without retrieval augmentation in terms of final \acc~and \cit~ in Table~\ref{table:result_multi}. 
OS-70B even matches or outperforms GPT4o on PubMedQA and QASA. 

\begin{table}[t!]
\centering
\small
    \resizebox{\textwidth}{!}{
    \begin{tabular}{l|cc|cc|cc||cc|cc|cc|c}
    \toprule
      & \multicolumn{6}{c||}{Single-paper performance} & \multicolumn{6}{c|}{Multi-paper performance} & Cost \\\midrule
      & \multicolumn{2}{c|}{Pub} & \multicolumn{2}{c|}{Sci} & \multicolumn{2}{c||}{QASA} & \multicolumn{2}{c|}{CS} & \multicolumn{2}{c|}{Multi} & Bio & Neu & CS \\ 
    \textbf{Model} & \acc & \cit & \acc & \cit & \acc & \cit &  \acc & \cit & \llm &  \cit & \cit & \cit  & USD / q \\  \midrule 
    Llama3-8B & 61.5 & 0.0 & 66.8 & 0.0 & 14.3 & 0.0 &  41.9 & 0.0 & 3.79 & 0.0 & 0.0 & 0.0 & 0.0001 \\ 
    +OSDS & 75.2 & 63.9 & 75.5 & 36.2 & 18.6 & 47.2 &   46.7 & 26.1 & \bf 4.22 & 25.3 & 38.0 & 36.8 & 0.0001 \\
    {\bf OS-8B} & {\bf 76.4} & {\bf 68.9} & {\bf 76.0} & {\bf 43.6} & \bf 23.0 & \bf 56.3   &  \bf 51.1 & \bf 47.9 &4.12 & \bf 42.8 & {\bf 50.8} & {\bf 56.8} & 0.003 \\\midrule
    Llama3-70B & 69.5 & 0.0 & 76.9 & 0.0 & 13.7 & 0.0   & 44.9 & 0.0 & 3.82 & 0.0 & 0.0 & 0.0 & 0.0004\\ 
    +OSDS & 77.4 & 71.1  & 78.2 & 42.5 & 22.7 & 63.6  &  48.5 & 24.5 & 4.24 & 41.4 &  53.8  & 58.1 & 0.0004 \\
    {\bf OS-70B} & \bf 79.6 & \bf 74.0 & {\bf 82.1} &\bf 47.5 & {\bf 23.4} & \bf 64.2 & {\bf 52.5} & \bf 45.9 &  4.03 & \bf 54.7  & {\bf 55.9} & {\bf 63.1} & 0.01 \\\midrule
    GPT4o & 65.8 & 0.0 & 77.8 & 0.0 & \bf 21.2 & 0.0   & 45.0 & 0.1 & 4.01 & 0.7 & 0.2 & 0.1  & 0.006  \\ 
    +OSDS & \bf 75.1 & 73.7 & 79.3 & 47.9 & 18.3 & 53.6 & 52.4 & 31.1 & 4.03  & 31.5 & 36.3  & 21.9 & 0.01 \\
{\bf OS-GPT4o} &  74.8 &  \bf 77.1  & \bf 81.3 &\bf 56.5   &18.7   &\bf 60.4 &{\bf  57.7} &  \bf 39.5  & {\bf 4.51} & \bf 37.5 & \bf{51.5} & \bf 43.5  & 0.05 \\\midrule
    PaperQA2 & -- & -- & -- & -- & -- & -- & 45.6 & {48.0} & 3.82 & { 47.2} & {56.7} & 56.0 & 0.3$\sim$2.3 \\
    Perplexity & -- & -- & -- & -- & -- & -- & 40.0 & -- & 4.15 & -- & -- & -- & 0.002$^{**}$\\
    \bottomrule
    \end{tabular}
    }
\caption{{\bf Results of \data}. CS, Multi, Bio and Neu indicate \textsc{Scholar-CS}, \textsc{ScholarQA-Multi}, \textsc{ScholarQA-Bio} and \textsc{ScholarQA-Neuro}, respectively. 
\acc~indicates correctness metrics (accuracy for PubMedQA and SciFact, ROUGE-L for QASA and Overall scores for \textsc{ScholarQA-CS}) and \cit~indicates citation F1. \llm~indicates the average score of \org~(organization), \rel~(relevance), \cov~(coverage) as predicted by Prometheus~\citep{kim2024prometheus}. $^*$PaperQA2 is based on GPT4o, and its pricing is dependent on the number of PDF files used during inference. 
For the 8B and 70B model costs, while evaluations were conducted on our local machines, we estimated costs based on Together.ai pricing.
$^{**}$We used Perplexity Pro (which requires a monthly subscription at \$20 USD) and divided this cost by 9,000, which is the maximum number of queries allowed under the Pro subscription. 
Since the Perplexity UI does not provide snippets for each citation, we were unable to evaluate its citation accuracy. 
}
\label{table:result_multi}
\end{table}%
\begin{table}
\small
    \centering
        \resizebox{\textwidth}{!}{
    \begin{tabular}{l| ccc|ccc}
    \toprule
        & \multicolumn{3}{c}{Computer Science} & \multicolumn{3}{c}{Biomedicine} \\ 
      Model  &Total \# & \# of Hallucinated  ($\downarrow)$ & Ratio ($\downarrow)$ &  Total \# &  \# of Hallucinated  ($\downarrow)$ & Ratio ($\downarrow)$  \\
    \midrule
     OS-8B &  9.65 &  0.0 & 0.0 & 6.25 &  0.0 & 0.0 \\ \hdashline
    Llama 3.1 8B &  5.20 & 4.79 & 92.1\% & 5.58 & 5.46 & 97.6\% \\
    Llama 3.1 70B & 6.14 & 4.78 & 78.1\% & 6.98 & 6.74 & 96.6\%\\
    GPT4o  & 5.74 & 4.52 & 78.7\% & 5.24  & 4.97 & 94.8\% \\
    % GPT4o & \checkmark & & & \\ 
    % % Llama3.1 70B  &  \checkmark & & & \\ 
    % Ours 8B &  \checkmark \\
    \bottomrule
    \end{tabular}
    }
\caption{{\bf Statistics of hallucinated papers in computer science and biomedicine domains}. Our analysis revealed a significant number of non-existent cited papers in predictions made by LLMs without retrieval, which is a problem not observed in \model.}\label{table:result_non_parametric}
\end{table}
\paragraph{Multi-paper tasks.} 
\model-8B, 70B, and GPT4o (OS-8B, OS-70B and OS-GPT4o) demonstrate strong performance in multi-paper tasks. Specifically, OS-GPT4o provides a 12.7 point improvement over GPT4o alone in \textsc{Scholar-CS} \acc~and a 5.3 improvement over standard RAG. 
When combined with trained OS-8B, \model significantly outperforms the pipeline that uses off-the-shelf Llama 3.1 8B, showcasing the benefits of domain-specific training. 
Furthermore, this \model-8B outperforms proprietary systems such as GPT4o, Perplexity Pro, or PaperQA2, which uses GPT4o models for passage reranking, summarization and answer generation, by a substantial margin. 
Notably, by leveraging efficient retrieval pipelines with lightweight bi-encoders, cross-encoders, and in-house models, \model-8B and \model-GPT4o achieve significantly lower costs—orders of magnitude cheaper than PaperQA2---while maintaining high performance.  

\paragraph{Limitations of parametric LMs.}
On both single-paper and multi-paper tasks, we observe that non-retrieval augmented baselines struggle and retrieval is almost always conducive to achieving better performance, and models without any retrieval often struggle to generate correct citations and show limited coverage on multi-paper tasks. 
As shown in Table~\ref{table:result_non_parametric}, the proportion of cited papers that actually exist is strikingly low. In particular, while models such as GPT4o and Llama can generate plausible reference lists, we find that 78-98\% of the cited papers are fabricated, and the issue is exacerbated in biomedical domains. Even when citations refer to real papers, the majority are not substantiated by the corresponding abstracts, resulting in near-zero citation accuracy. 

We also observe that such models also generate responses with limited coverage. 
On \textsc{Scholar}-Multi, non-retrieval models (Llama 3.1 8B, 70B, and GPT4o) consistently exhibit significantly lower average scores compared to retrieval-augmented models. This discrepancy is largely driven by much lower \cov~scores; for instance, Llama 3.1 8B achieves a \cov~score of 3.45, while Llama 3.1 8B + OSDS (a standard RAG baseline) improves it to 4.01. 
These results suggest that solely relying on models' parametric knowledge alone is particularly difficult in scientific domains, especially for smaller LMs. 

\subsection{Analysis}
\begin{figure}[t!]
\begin{subfigure}[b]{0.3\textwidth}
\centering
\footnotesize
\begin{tabular}{lcc}
\toprule
 & \multicolumn{2}{c}{\textsc{Scholar-CS}}  \\
 & \acc & \cit \\
\midrule
\textsc{OS}-8B & 51.3 & 47.9  \\ \hdashline
- training & 49.4 &  42.3 \\
- reranking & 49.6 & 28.2  \\
% -  meta filter & 49.2 & 36.7 \\
- feedback & 51.1 &  50.2  \\
-  attribution & 49.3 & 44.0  \\
\midrule
OS-\textsc{GPT4o} & 57.7  & 39.5 \\ \hdashline
-  reranking &  52.4  & 22.9 \\
-  feedback & 55.1 & 31.0  \\
-  attribution & 55.6 & 30.6  \\
\bottomrule
\end{tabular}
\caption{Ablation of different components of \model.}
\label{tab:ablatioon}
\end{subfigure}%
  \hspace{1.2cm}
  \begin{subfigure}[b]{0.28\textwidth}
        \includegraphics[width=\textwidth]{figs/top_n_accuracy.pdf}
        \caption{Top N Ablations (\acc) on \textsc{Scholar-CS}. }
        \label{fig:top_n_ablations_accuracy}
  \end{subfigure}%
  \hspace{0.3cm}
  \begin{subfigure}[b]{0.27\textwidth}
        \includegraphics[width=\textwidth]{figs/top_n_citation.pdf}
        \caption{Top N Ablations (\cit-F1) on \textsc{Scholar-CS}. }
        \label{fig:top_n_ablations_citation}
  \end{subfigure}%
\caption{{\bf Analysis on \model:} (a) {\bf Ablation studies} for key components of \model training and inference based on different underlying LMs.  (b) {\bf Top N docs:} Analysis of the effect of varying the number of context chunks for final downstream tasks. We evaluate final model performance based on citation accuracy and correctness on multi-doc QA tasks, using \model 8B and Llama 3.1 8B. }\label{fig:analysis}
\end{figure}

\paragraph{Ablation studies.}
We conduct ablations to assess the effectiveness of individual components of \model (inference and training). Specifically, we remove each of the inference-time procedures, reranking, feedback and attribution, and for OS-8B, we ablate the training, where we use Llama3-8B without any further training.  

As shown in Figure~\ref{fig:analysis} (a), removing these components significantly impacts both the overall correctness and citation accuracy of the model outputs. Notably, removing the reranker led to substantial performance drops on both models. 
The pronounced decline in performance after removing feedback loops in  GPT4o suggests that more powerful models greatly benefit from a self-feedback cycle, consistent with ~\citet{madaan2023self}, while it gives limited performance drops in our trained 8B.  Additionally, the removal of post-hoc attribution assessments negatively affected both citation accuracy and final output correctness, highlighting the importance of ensuring that models verify their outputs. 
The significant performance gap between trained versus vanilla OS-8B suggests that further training on high-quality, domain-specific data is key to building efficient, task-specialized LMs. 
In the next analysis, we demonstrate that training has a significant impact on an LM's ability to effectively utilize more context, while maintaining citation accuracy. 

\paragraph{Number of context passages.}
We analyzed how varying the number of context passages  (top $N$) impacts model performance. 
Specifically, we experimented with Standard \texttt{RAG} and \model using our trained 8B model and Llama 3.1 8B, and evaluated both generation accuracy and citation accuracy on \textsc{Scholar-CS}. 
Figures~\ref{fig:analysis} (b)(c) show the results.
Although Llama 3.1 is trained to handle and accept a context length of up to 128K tokens, we found that its performance deteriorates after a certain context size. While increasing the top $N$ context window from 5 to 10 does improve the model’s correctness score, expanding further actually worsens both correctness and citation accuracy. This suggests that even though LMs can process large numbers of passages, they may struggle to effectively use them without specialized training, particularly for smaller models. 

In contrast, our trained 8B model maintains strong performance for up to $N=20$ passages. 
We also found larger models such as Llama 3.1 70B to be more robust to increased context length. 
In terms of citation accuracy, as shown in Figures~\ref{fig:analysis} (c), Llama 3.1 8B observes quick decline and citation F1 gets as low as 10, while our 8B LM and Llama 70B both maintain around 40 citation F1, although they also see some performance decline. 